% p.29 — الباب الخامس العشر في رمى يرمي (تابع): المنفيات والمؤكدات
\mspage{29}
% Layout: six tall columns (right-to-left), each a stack of boxed blocks of unequal height.
% Every sub-heading box holds the heading AND, beneath it, the underlying (pre-iʿlāl) form,
% e.g. «مجهول بلم / لَمْ يُرْمَيْ»; the regular forms follow outside the box. Row alignment is
% approximate: all six columns run to 21 lines, so rows are aligned 1:1; the \cline marks
% reproduce the printed box boundaries inside each column.
% Column 1 continues the مضارع بلم (معروف) block begun on p.28. In the ناقص يائي the underlying
% and the actual forms of the مؤكدات (e.g. لَا يَرْمِيَنَّ) coincide, so the lithograph prints them twice.
% The running header line of the frame is blank (smudged) on this page.
\begin{sarf}{|c|c|c|c|c|c|}
لَمْ تَرْمِيَا & لَمْ أُرْمَ & لَنْ يُرْمَيَا & لَتَرْمِيَنَّ & لَتُرْمَيْنَانِّ & \thead{أمر حاضر بالنون}\\
لَمْ تَرْمُوْا & لَمْ نُرْمَ & لَنْ يُرْمَوْا & لَتَرْمِيَانِّ & لَأُرْمَيَنَّ & اِرْمِيَنَّ\\\cline{2-2}\cline{6-6}
لَمْ تَرْمِيْ & \thead{مضارع بلن} & لَنْ تُرْمَى & لَتَرْمُنَّ & لَنُرْمَيَنَّ & اِرْمِيَنَّ\\\cline{5-5}
لَمْ تَرْمِيَا & لَنْ يَرْمِيَ & لَنْ تُرْمَيَا & لَتَرْمِنَّ & \thead{منفي بالنون} & اِرْمِيَانِّ\\\cline{2-2}
لَمْ تَرْمِيْنَ & لَنْ يَرْمِيَ & لَنْ يُرْمَيْنَ & لَتَرْمِيَانِّ & لَا يَرْمِيَنَّ & اِرْمُنَّ\\\cline{5-5}
لَمْ أَرْمِ & لَنْ يَرْمِيَا & لَنْ تُرْمَى & لَتَرْمِيْنَانِّ & لَا يَرْمِيَنَّ & اِرْمِنَّ\\\cline{5-5}
لَمْ نَرْمِ & لَنْ يَرْمُوْا & لَنْ تُرْمَيَا & \com{لَاَرْمِيَنَّ}{كذا في الأصل، بألف بلا همزة؛ ولعل الصواب: لَأَرْمِيَنَّ} & \thead{مجهول بالنون} & اِرْمِيَانِّ\\\cline{1-1}
\thead{مجهول بلم} & لَنْ تَرْمِيَ & لَنْ تُرْمَوْا & لَنَرْمِيَنَّ & لَا يُرْمَيَنَّ & اِرْمِيْنَانِّ\\\cline{4-4}\cline{5-5}\cline{6-6}
لَمْ يُرْمَيْ & لَنْ تَرْمِيَا & لَنْ تُرْمَيْ & \thead{مجهول بالنون} & لَا يُرْمَيَنَّ & \thead{نهي حاضر بالنون}\\\cline{1-1}\cline{5-5}
لَمْ يُرْمَ & لَنْ يَرْمِيْنَ & لَنْ تُرْمَيَا & لَيُرْمَيَنَّ & \thead{أمر غائب بالنون} & لَا تَرْمِيَنَّ\\\cline{4-4}\cline{6-6}
لَمْ يُرْمَيَا & لَنْ تَرْمِيَ & لَنْ تُرْمَيْنَ & لَيُرْمَيَنَّ & لِيَرْمِيَنَّ & لَا تَرْمِيَنَّ\\\cline{5-5}
لَمْ يُرْمَوْا & لَنْ تَرْمِيَا & لَنْ أُرْمَى & لَيُرْمَيَانِّ & لِيَرْمِيَنَّ & لَا تَرْمِيَانِّ\\\cline{5-5}
لَمْ تُرْمَ & لَنْ تَرْمُوْا & لَنْ نُرْمَى & لَيُرْمَوُنَّ & \thead{مجهول بالنون} & لَا تَرْمُنَّ\\\cline{3-3}
لَمْ تُرْمَيَا & لَنْ تَرْمِيْ & \thead{مضارع بالنون} & لَتُرْمَيَنَّ & لِيُرْمَيَنَّ & لَا تَرْمِنَّ\\\cline{5-5}
لَمْ يُرْمَيْنَ & لَنْ تَرْمِيَا & لَيَرْمِيَنَّ & لَتُرْمَيَانِّ & لِيُرْمَيَنَّ & لَا تَرْمِيَانِّ\\\cline{3-3}\cline{5-5}
لَمْ تُرْمَ & لَنْ تَرْمِيْنَ & لَيَرْمِيَنَّ & لَيُرْمَيْنَانِّ & \thead{نهي غائب بالنون} & لَا تَرْمِيْنَانِّ\\\cline{6-6}
لَمْ تُرْمَيَا & لَنْ أَرْمِيَ & لَيَرْمِيَانِّ & لَتُرْمَيَنَّ & لَا يَرْمِيَنَّ & \thead{تمرين}\\\cline{5-5}\cline{6-6}
لَمْ تُرْمَوْا & لَنْ نَرْمِيَ & لَيَرْمُنَّ & لَتُرْمَيَانِّ & لَا يَرْمِيَنَّ & هَدٰى\\\cline{2-2}\cline{5-5}
لَمْ تُرْمَيْ & \thead{مجهول بلن} & لَتَرْمِيَنَّ & لَتُرْمَوُنَّ & \thead{مجهول بالنون} & جَنٰى\\
لَمْ تُرْمَيَا & لَنْ يُرْمَيَ & لَتَرْمِيَانِّ & لَتُرْمَيِنَّ & لَا يُرْمَيَنَّ & كَفٰى\\\cline{2-2}\cline{5-5}
لَمْ تُرْمَيْنَ & لَنْ يُرْمَى & لَيَرْمِيْنَانِّ & لَتُرْمَيَانِّ & لَا يُرْمَيَنَّ & مَشٰى\\
\end{sarf}
